

Tessera's gains over expression-only clustering in lung tissue are real but are not evidence for its distinctive components. A neighbor mean whose scale was chosen on other patients outperformed it in every bounded lung input. Whole-section sensitivity analyses retained this ordering with a fixed 200-neighbor control informed by the same cohort; they test the sampling budget, not independent generalization. The lung regions were constructed from cell-type frequencies among 200 neighbors, and averaging at that scale recovers them without learning. An independent cohort whose regions a pathologist drew on the stained tissue, and which therefore was not built by neighborhood averaging, gave the same ordering: an untuned mean on Tessera's own graph exceeded the encoder in most patients, and so did a control tuned on the other patients. Spatial encoders should face neighborhood baselines with a comparable opportunity to choose their scale, and an unweighted same-graph mean, which needs neither labels nor tuning, is the cheapest informative comparator.

The component tests did not show a useful contribution from learned gating to lung agreement. The gates had small magnitudes, separated seams less well than a plain expression difference, and replacing them by uniform weights caused no systematic loss across these patients. Their relative weights were still heterogeneous after normalization, so small magnitude does not make the operator an unweighted mean. The ablation provides the mechanistic evidence: message passing with reconstruction alone reproduced the self-only encoder, while the neighbor-positive contrastive objective supplied the average lung gain. This finding concerns this released implementation, not gated graph networks in general; different gate objectives would need their own evaluation on untouched patients. In breast, gates gave near-chance tumor-seam AUROC; three anatomical tissues gave below-chance values, within the earlier 24-object median of \LandGateMedian{}. In colorectal tissue they carried weak seam information without demonstrated added value. Smoothing the feature difference over a neighborhood gave a better seam score, whereas reweighting a six-neighbor average by feature similarity showed no consistent advantage, so expression change at the seams carries information that neither the learned gate nor this simple reweighting turned into better domains at this resolution.

Target and clustering resolution both matter. At the primary cluster number, molecular or dominant-category targets favored expression-only clustering on BCC, breast, Open-ST and MIBI. Reducing the cluster number made BCC mixed, so that negative ordering is not universal. Within lung, the same embeddings that gained on regions lost on cell types, and transfer reproduced this trade-off with a different assignment rule. The 200-neighbor transfer control was informed by within-object analyses and is descriptive, although its transforms remain source-only. BANKSY explicitly weights the cell and its neighborhood \citep{singhal2024banksy}; our lightly weighted BANKSY-style control, not a native BANKSY run, was best in the bounded BCC selection. Multi-scale and composition representations address the region side of this trade-off \citep{yuan2024mender,liang2026cellniche}. Tessera's advantage was not uniform across annotations. It exceeded the same-graph mean in smooth anatomical tissue and in most computed lung regions but rarely in the pathologist-labelled patients, and it exceeded the best simple control chosen with the labels only in the cortex section used in its development and in colorectal TVA. Those exceptions show where a learned component could add structure beyond a neighborhood average, but one development section cannot establish it.

These results bear on how spatial methods are benchmarked. Comparisons with expression-only baselines reward any smoothing on spatially constructed targets, and untuned simple baselines understate what position alone can achieve. Benchmarks that already compare many native methods across platforms \citep{yuan2024benchmark,kang2025nar,chen2025imeta} would gain from including tuned and same-graph neighborhood averages and from reporting cell-type and domain targets separately. Seed variation also deserves attention: Tessera's partitions changed between seeds far more than PCA's, and repeated fits of identical code changed per-patient scores by up to \RepMaxRange{} ARI, so a single fitted partition is a weak basis for interpreting domain boundaries or for ranking methods whose scores differ by less than that.

Several limits apply. The study evaluates one released encoder at fixed hyperparameters; tuning it could change the findings, and fine-tuned or other spatial encoders were not tested. Five lung, four BCC and eight breast patients give coarse sign-flip summaries, and overlapping selection folds limit their inferential interpretation. The independent cohort has one pathologist, no second annotator and image registration that we inherited, and its weighting controls, boundary variants and tissue panel were chosen after the first comparison, so they are exploratory. The three non-tumor tissues are single sections that the encoder was benchmarked on earlier. Other annotations are computational or transferred labels, and the cluster number is label-informed. The Open-ST analysis remains a thinned subset, not a whole-section validation, and registered histology and protein channels support inspection only (Supplementary Fig.~S3). The three-patient colorectal pilot adds an untouched source, but cost selection, mixed labels and its disclosed QC sensitivity limit confirmation. Author H\&E overlays are not independent registration review. Stronger validation needs blinded adjudicated annotations and a patient sample designed for an effect margin. Calibration and selective abstention would require a separate protocol if uncertainty were claimed. A decisive test would use boundaries annotated independently by at least two readers in patients that informed neither method nor scale selection, together with a calibrated uncertainty readout.

In the tested lung and independent breast cohorts, simple neighborhood averaging accounts for much of the apparent gain over expression-only clustering. Component ablations identify the neighbor-positive contrastive objective as the source of the lung gain, while learned gates were dispensable there. Matched and tuned neighborhood averages should accompany architectural comparisons, with target construction, patient separation and computational variation reported explicitly.

